\chapter{用户研究与需求分析}\label{chap:requirement}

任何看似显然的设计判断都应当先经得起一次「我们到底在为谁解决什么问题」的追问。本章的用户研究就服务于这一追问。受课程周期所限，本研究并未追求统计意义上的代表性，但仍尽可能在「量化 + 质化 + 群体 + 个体」四个维度上互相印证，从中提炼出 Mood Radio 的目标用户、需求边界与设计原则。

\section{研究方法}

为了让结论不只是建立在一份问卷之上，课题组采用了四种相互补充的方法。问卷面向较大样本，用以采集行为频率与态度倾向；情景访谈则借由「就地观察 + 半结构化提问」深入用户实际的使用环境，去捕捉那些用户自己未必能在问卷里讲清楚的微小动作；焦点小组用于让不同立场的同学相互辩论，逼出关于「系统是否应该改变用户情绪」、「群体在场是否舒适」等带价值判断的问题；最后，单独访谈则用于深挖那些在问卷里表达出较强隐私顾虑的少数案例，理解他们的边界是如何被划出的。四种方法的样本量与覆盖目的汇总在\cref{tab:research-methods}~中。

\begin{table}[!htb]
  \centering
  \caption{用户研究方法概览}\label{tab:research-methods}
  \begin{tabularx}{\linewidth}{lcX}
    \toprule
    \textbf{方法} & \textbf{样本量} & \textbf{覆盖目的} \\
    \midrule
    问卷调查    & 42 人 & 听歌习惯、情绪调节频率、找歌成本与隐私态度的量化采集 \\
    情景访谈    & 6 人  & 在用户真实使用环境下观察其「为什么打开音乐 App」 \\
    焦点小组    & 5 人  & 群体讨论「匹配 vs 安抚」与群体在场感的态度边界 \\
    单独访谈    & 4 人  & 深挖个体的失败案例与隐私担忧的具体边界 \\
    \bottomrule
  \end{tabularx}
\end{table}

\section{从问卷到访谈：研究发现}

问卷部分把 42 份回收答卷压缩成几个可比对的数字。\cref{tab:survey}~列出其中最为关键的四个题项。可以看到，绝大多数受访者会主动按心情挑歌（81\%），却又经常陷入「想听点什么但说不上来」的尴尬（76\%）。在隐私态度上，超过七成的人对「将面部或情绪数据上传到云端」持保留意见，而又有同样接近的比例希望系统能在「匹配心情」与「安抚心情」之间提供切换。这些数字一方面印证了\cref{chap:introduction}~中提出的问题确实是个真实痛点，另一方面也勾画出本系统不能轻易越过的两条边界：本地优先与策略可选。

\begin{table}[!htb]
  \centering
  \caption{问卷调查中四个关键题项的反馈比例（$N=42$）}\label{tab:survey}
  \begin{tabular*}{\linewidth}{@{\extracolsep{\fill}}lc}
    \toprule
    \textbf{题项} & \textbf{占比} \\
    \midrule
    会根据当下心情主动选择音乐 & 81\,\% \\
    经常出现「想听歌但不知道听什么」 & 76\,\% \\
    对将面部/情绪数据上传到云端存在顾虑 & 74\,\% \\
    希望系统能同时支持「匹配心情」与「安抚心情」两种策略 & 72\,\% \\
    \bottomrule
  \end{tabular*}
\end{table}

数字之外的细节由访谈补足。情景访谈中有一位受访者把场景描述得颇为生动——她在自习室坐了三个小时后准备放空，习惯性打开音乐，发现自己「就是不想点」，最终干脆什么都没听。这类「不想动」的瞬间几乎在每一位情景访谈对象身上都出现过，它意味着系统应当为「低操作」场景留出无障碍的入口；任何还要求用户输入关键词的设计在这一刻都会败下阵来。焦点小组中另一位同学讲了一句被课题组反复引用的原话：「我难过的时候不想被一首假装欢乐的歌打断。」这一句话直接否定了「以欢乐对抗低落」的简单粗暴策略，并促使我们在后续设计中把「匹配心情」放在与「安抚心情」同等重要的位置。至于单独访谈中的隐私担忧则相对集中，几位受访者把边界划在了「我能不能离线运行」这条线上——只要不联网，他们就愿意打开摄像头；一旦数据上行，他们就会立刻关掉。这一边界后来落到了\cref{tab:nfr}~中的 NFR-01 与 NFR-02 两条非功能性需求上。

\section{目标用户画像}

把研究发现凝聚成「想象中的人」是设计推进的有效手段。课题组从四种方法的样本里抽取出三个 Persona，它们在「是否愿意开摄像头」、「是否在意视觉沉浸」、「是否关心隐私」这三个维度上覆盖了相对极端但仍然合理的组合，目的是在后续设计取舍中始终有可参照的具体对象。三位 Persona 的画像如\cref{tab:persona}~所示。

\begin{table}[!htb]
  \centering
  \caption{Mood Radio 三个核心目标用户画像}\label{tab:persona}
  \begin{tabularx}{\linewidth}{>{\bfseries}l>{\raggedright\arraybackslash}p{3cm}X}
    \toprule
    \textbf{Persona} & \textbf{身份与背景} & \textbf{典型场景与核心诉求} \\
    \midrule
    小南（深夜情绪派）
      & 21 岁，本科生，理工科 &
      深夜赶完代码、心情起伏大，希望「一打开就有合适的歌」；不愿打字搜索，但乐于开摄像头。\\
    \addlinespace
    Aria（隐私敏感型）
      & 28 岁，独立设计师 &
      在咖啡馆办公，担心摄像头与情绪上传；坚持「电脑里所有数据都属于自己」；愿意手动选择情绪。\\
    \addlinespace
    阿陈（沉浸聆听者）
      & 24 岁，硕士在读，黑胶爱好者 &
      重视「视觉 + 听觉 + 物理感」的整体氛围；希望封面、转盘、背景跟着歌曲走；可以接受较多视觉变化。\\
    \bottomrule
  \end{tabularx}
\end{table}

三位 Persona 的差异性恰恰是「\emph{多通道冗余}」这一设计选择的根据：小南乐于使用摄像头，那么摄像头通道必须高效；Aria 拒绝摄像头，那么手动通道必须同等好用；阿陈愿意花时间在情绪图谱上探索，那么二维 V/A 图就需要被认真地做出来。三条通道在系统内部最终都会归并到同一个推荐入口，但在前端体验上始终保持各自独立的高完成度。

\section{需求与设计原则}

通过对前述发现的提炼，本系统面向的核心用户诉求可以概括为四点。\emph{首先}，用户希望「快」——在状态不佳时，最不希望被复杂的操作流程绊住。\emph{其次}，用户希望「自主」——是否「匹配」当下情绪、还是「安抚」当下情绪，应当由自己说了算，而非系统替自己决定。\emph{再次}，用户希望「看得懂」——系统至少要告诉自己「为什么是这首歌」，哪怕只是一个小小的徽标。\emph{最后}，用户希望「私密」——摄像头与音乐文件不应上传，而群体页面也要做匿名化处理。把这四点诉求落到工程语言上，就得到了\cref{tab:fr}（功能性需求）与\cref{tab:nfr}（非功能性需求）所列出的需求说明书。其中 H/M/L 三档优先级用于在课程时间内做合理的功能取舍。

\begin{table}[!htb]
  \centering
  \caption{Mood Radio 功能性需求（节选）}\label{tab:fr}
  \begin{tabularx}{\linewidth}{>{\ttfamily}l>{\bfseries}cX}
    \toprule
    \textbf{编号} & \textbf{优先级} & \textbf{描述} \\
    \midrule
    FR-01 & H & 系统应支持扫描本地音乐目录，并对每首歌计算 V/A 情绪坐标 \\
    FR-02 & H & 用户应能通过摄像头实时识别面部情绪 \\
    FR-03 & H & 用户应能手动从 7 类情绪标签中选择当前情绪 \\
    FR-04 & H & 用户应能通过自然语言聊天表达情绪并得到推荐 \\
    FR-05 & M & 用户应能基于当前时段一键触发推荐 \\
    FR-06 & M & 用户应能通过头部手势完成切歌等操作 \\
    FR-07 & H & 系统应在 SQL 层完成基于情绪距离、用户反馈、新鲜度与时间惩罚的复合打分推荐 \\
    FR-08 & M & 系统应能区分「匹配」与「安抚」两种策略并按此调整目标坐标 \\
    FR-09 & H & 系统应记录用户的播放历史与喜欢/不喜欢反馈 \\
    FR-10 & M & 系统应提供匿名群体在场感页面（情绪分布、Top 听单） \\
    FR-11 & M & 系统应记录用户的心情日志并以日历/堆叠图可视化 \\
    FR-12 & L & 系统应提供可配置的最近播放排除时长、推荐半径等偏好参数 \\
    \bottomrule
  \end{tabularx}
\end{table}

\begin{table}[!htb]
  \centering
  \caption{Mood Radio 非功能性需求（节选）}\label{tab:nfr}
  \begin{tabularx}{\linewidth}{>{\ttfamily}l>{\bfseries}cX}
    \toprule
    \textbf{编号} & \textbf{优先级} & \textbf{描述} \\
    \midrule
    NFR-01 & H & 摄像头画面与本地音乐文件均\emph{不}上传至云端 \\
    NFR-02 & H & LLM 的 API Key 不写入数据库或后端日志，仅存放于浏览器 \\
    NFR-03 & H & 推荐响应时间 $\leq 500$\,ms（曲库 $\leq 200$ 首时） \\
    NFR-04 & M & 系统支持完全离线运行（聊天功能除外） \\
    NFR-05 & M & 系统视觉应随当前情绪柔和变化，过渡时长 $\geq 400$\,ms 以避免突变 \\
    NFR-06 & M & 群体在场页面的隐私门槛可配置（默认 $N=1$，可设为 $\geq 3$） \\
    NFR-07 & L & 主要交互应可通过键盘完成（Skip link、aria 标签） \\
    \bottomrule
  \end{tabularx}
\end{table}

\section{需求框图}

把上述功能性与非功能性需求按模块组织起来，便得到了\cref{fig:req-diagram}~所示的需求框图。它以「情绪感知音乐电台」为根，自然地派生出情绪输入、音乐情绪建模、推荐与反馈、视觉沉浸、群体在场与隐私保护六个一级模块。这六个模块不仅在功能上彼此独立，更重要的是它们对应了不同的设计权衡：情绪输入需要在准确度与噪声鲁棒性之间取舍，音乐情绪建模需要在精度与计算成本之间取舍，推荐与反馈需要在距离与新鲜度之间取舍，而隐私保护则要求始终对前几个模块做出抑制。

\begin{figure}[!htb]
  \centering
  \begin{tikzpicture}[
      every node/.style={font=\small},
      root/.style={draw, rounded corners=2pt, align=center,
        inner sep=4pt, fill=orange!15, font=\small\bfseries},
      mod/.style={draw, rounded corners=2pt, align=left,
        inner sep=3pt, fill=blue!8, text width=3.2cm, minimum height=0.7cm},
      leaf/.style={draw, rounded corners=2pt, align=left,
        inner sep=3pt, fill=gray!8, font=\scriptsize, text width=3.0cm, minimum height=0.55cm},
    ]
    \node[root] (r) at (0, 0) {Mood Radio \\ 情绪感知音乐电台};

    \node[mod] (m1) at (4.5,  3.2) {情绪输入 \\ \scriptsize FR-02$\sim$06};
    \node[mod] (m2) at (4.5,  1.8) {音乐情绪建模 \\ \scriptsize FR-01};
    \node[mod] (m3) at (4.5,  0.4) {推荐与反馈 \\ \scriptsize FR-07/08/09};
    \node[mod] (m4) at (4.5, -1.0) {视觉沉浸 \\ \scriptsize NFR-05};
    \node[mod] (m5) at (4.5, -2.4) {群体在场 \\ \scriptsize FR-10, NFR-06};
    \node[mod] (m6) at (4.5, -3.8) {隐私保护 \\ \scriptsize NFR-01/02};

    \node[leaf] (l1) at (9.5,  4.4) {摄像头 / 手动};
    \node[leaf] (l2) at (9.5,  3.4) {聊天 / 时段 / 手势};
    \node[leaf] (l3) at (9.5,  1.8) {扫描 + V/A 标注};
    \node[leaf] (l4) at (9.5,  0.4) {复合打分 + 反馈};
    \node[leaf] (l5) at (9.5, -1.0) {主题 / Aurora / 黑胶};
    \node[leaf] (l6) at (9.5, -2.4) {匿名身份 + 聚合};
    \node[leaf] (l7) at (9.5, -3.8) {本地推理 + Key 不落盘};

    \foreach \dst in {m1, m2, m3, m4, m5, m6} \draw[-, thick] (r) -- (\dst);
    \draw[-] (m1) -- (l1); \draw[-] (m1) -- (l2);
    \draw[-] (m2) -- (l3); \draw[-] (m3) -- (l4);
    \draw[-] (m4) -- (l5); \draw[-] (m5) -- (l6);
    \draw[-] (m6) -- (l7);
  \end{tikzpicture}
  \caption{需求框图：以六个模块组织 FR/NFR}\label{fig:req-diagram}
\end{figure}

\section{设计原则}

最终落到设计上，本系统给自己定下了五条贯穿全局的原则。「\textbf{低负担输入}」要求系统始终给用户保留一种「一眼一击」的最快通道，并通过多种模态的冗余避免任何单一传感器的失效拖累整体；「\textbf{用户可控}」把匹配/安抚、自动续推、摄像头开关、最近播放排除时长等关键开关放在显著位置，使「让系统替我做主」始终是用户的主动选择而非默认结果；「\textbf{可解释}」要求每一次推荐都能通过情绪图谱、推荐徽标或聊天解释告诉用户「为什么是这首」；「\textbf{本地优先}」则贯彻了用户对隐私的诉求——摄像头与音乐均本地处理，LLM 密钥只存放于浏览器，群体页面采用聚合统计；最后，「\textbf{沉浸感与一致性}」要求视觉、声音、文字三个层面共同响应同一个情绪标签，避免不同通道之间互相打架。这五条原则在表面上略显抽象，但它们将以非常具体的方式出现在后续章节的每一个设计判断中。值得一提的是，「沉浸感」并不是出于审美追求；让用户在视觉层面持续感受到系统对其情绪的「跟随」，能够显著减少他们对识别准确度的焦虑，进而提高对推荐结果的接受度，这一点与文献中关于反馈一致性对用户信任影响的讨论是一致的\cite{shneiderman2016designing}。
